\documentclass[11pt]{article}

\usepackage[T1]{fontenc}
\usepackage{amsmath,amssymb,amsthm}
\usepackage{hyperref}

\title{A Regular Asymptotic Roadmap for SEM Outputs}
\author{}
\date{\today}

\newcommand{\E}{\mathbb{E}}
\newcommand{\Var}{\operatorname{Var}}
\newcommand{\Cov}{\operatorname{Cov}}
\newcommand{\tr}{\operatorname{tr}}
\newcommand{\argmin}{\operatorname*{arg\,min}}
\newcommand{\dto}{\xrightarrow{d}}
\newcommand{\pto}{\xrightarrow{p}}

\theoremstyle{plain}
\newtheorem{theorem}{Theorem}
\newtheorem{corollary}{Corollary}
\theoremstyle{remark}
\newtheorem*{remark}{Remark}

\begin{document}
\maketitle

\section*{Purpose}

This note sketches a paper and implementation roadmap for regular
misspecification asymptotics in structural equation modeling. The aim is a
math-forward statistical-methods paper, paired with implementation in
\texttt{magmaan}, that treats many familiar SEM quantities as consequences of a
small number of profile-map theorems.

The scope is deliberately regular:
\[
  \text{interior pseudo-true points, fixed rank, smooth maps, nonsingular
  profile Hessians.}
\]
Boundary parameters, rank loss, unidentified pseudo-true points, chi-bar-square
limits, clipping, and other nonregular phenomena should be acknowledged, but
not solved in the first paper. They are too large to include without blurring
the central message.

The main claim is not that every mathematical ingredient is new. The claim is
that the ingredients are scattered, repeatedly rederived, and often
implemented through exact-fit shortcuts. A unified regular theory can make clear
which object is being estimated, which Hessian is required under
misspecification, and how standard errors, fit measures, nested tests,
relative-fit tests, and future indices such as CFI and TLI share the same
calculus.

\section{Paper Positioning}

The paper should be pitched as a statistical-methods contribution with broad SEM
impact. A JASA-style target is plausible because the procedures are applied at
large scale, the asymptotic objects are nontrivial, and a software-backed
unification would be useful beyond one psychometric subcommunity.

The paper should not argue that the work is ``too mathematical'' for
psychometrics journals. That is unnecessary and would read badly. A better
framing is:
\[
\begin{minipage}{0.86\linewidth}
Routine SEM output is built from profile estimators and profiled discrepancy
functions. Under misspecification, the relevant derivatives are observed
profile derivatives at pseudo-true points, not the exact-fit/Fisher
simplifications often used in derivations and software. We provide a unified
regular asymptotic theory and an implementation that applies the same objects to
standard errors, model-comparison tests, robust RMSEA, and related fit indices.
\end{minipage}
\]
This is candid about the contribution: not a new probability limit theorem for
its own sake, but a unifying theorem, a correction of common shortcuts, and an
implementation-ready framework for procedures that are used constantly.

\section{Master Setup}

Let $\hat x$ be the first-stage object used by the SEM estimator:
sample means/covariances, FIML score aggregates, thresholds and polychorics,
group proportions, NACOV components, estimated weights, or any stacked
combination needed to make later maps smooth. Assume
\begin{equation}\label{eq:first-stage}
  \hat x
  =
  x_0+n^{-1/2}Z+n^{-1}b_n+o_p(n^{-1}),
  \qquad
  Z\dto N(0,\Gamma),
  \qquad
  \E b_n\to b .
\end{equation}
The $n^{-1}b$ term is optional for first-order inference but needed for
second-order bias corrections such as fixed-misspecification RMSEA.

For each model $M$, define a smooth criterion
\[
  f_M(\theta,x),
  \qquad
  \theta_M(x)=\argmin_\theta f_M(\theta,x),
  \qquad
  v_M(x)=f_M\{\theta_M(x),x\}.
\]
Here $v_M$ is the profiled discrepancy or value function. In a quadratic
moment-discrepancy setting,
\[
  f_M(\theta,x)
  =
  \frac12 r_M(\theta,x)'W(x)r_M(\theta,x),
\]
but the theory should be written for a general smooth criterion so that ML,
FIML, GLS, ULS, DWLS, ordinal limited-information estimators, and multigroup
extensions fit the same template.

All derivatives below are evaluated at $(\theta_M,x_0)$, where
$\theta_M=\theta_M(x_0)$ is the pseudo-true parameter. Write
\[
  H_M=f_{\theta\theta},
  \qquad
  G_M=f_{\theta x}.
\]
The regularity assumption is that $H_M$ is nonsingular after applying any fixed
linear equality constraints or reduced parameterization.

\section{Main Theorem 1: Coefficient Maps}

The first master map is the pseudo-true parameter map $x\mapsto\theta_M(x)$.
The population first-order condition is
\[
  f_\theta\{\theta_M(x),x\}=0 .
\]
Differentiating gives
\[
  D\theta_M[x_0] = -H_M^{-1}G_M .
\]
Therefore
\begin{equation}\label{eq:theta-if}
  \sqrt n(\hat\theta_M-\theta_M)
  =
  -H_M^{-1}G_M Z+o_p(1),
\end{equation}
and
\begin{equation}\label{eq:theta-vcov}
  \Var\{\sqrt n(\hat\theta_M-\theta_M)\}
  =
  H_M^{-1}G_M\Gamma G_M'H_M^{-1}.
\end{equation}

This is the master theorem for misspecification-robust standard errors. In SEM
notation, $H_M$ is the observed/profile Hessian at the pseudo-true point. Under
correct specification it may reduce to expected information or a
Gauss--Newton/Fisher expression, but that reduction is a special case, not the
general theorem.

\section{Main Theorem 2: Residual and Fitted-Value Maps}

Many SEM outputs are not just parameters or scalar discrepancies. They use the
fitted moments, residuals, standardized residuals, modification indices, EPCs,
or score-like projections. These should be handled as smooth maps derived from
$\theta_M(x)$.

For a residual map
\[
  r_M(x)=r_M\{\theta_M(x),x\},
\]
the derivative is
\begin{equation}\label{eq:resid-derivative}
  Dr_M
  =
  r_x-r_\theta H_M^{-1}G_M ,
\end{equation}
with signs determined by the residual convention. Thus
\[
  \sqrt n\{\hat r_M-r_M\}
  =
  Dr_M\,Z+o_p(1).
\]
This is the natural home for residual diagnostics, SRMR ingredients,
standardized residuals, score-test ingredients, and modification-index
building blocks. It is not a new asymptotic primitive, but it deserves a named
corollary because many SEM procedures are projection procedures rather than
coefficient procedures.

\section{Main Theorem 3: Profiled Discrepancy Maps}

The second master primitive is the profiled value map
\[
  v_M(x)=\min_\theta f_M(\theta,x).
\]
The envelope theorem gives
\[
  s_M=v_{M,x}=f_x ,
\]
and the profile Hessian is
\begin{equation}\label{eq:Q}
  Q_M
  =
  v_{M,xx}
  =
  f_{xx}-f_{x\theta}H_M^{-1}f_{\theta x}.
\end{equation}
Then
\begin{equation}\label{eq:value-expansion}
  v_M(\hat x)
  =
  v_M(x_0)
  +n^{-1/2}s_M'Z
  +n^{-1}\left\{s_M'b_n+\frac12 Z'Q_MZ\right\}
  +o_p(n^{-1}).
\end{equation}
Consequently
\begin{equation}\label{eq:value-bias}
  \E\{v_M(\hat x)\}
  =
  v_M(x_0)
  +n^{-1}\left\{s_M'b+\frac12\tr(Q_M\Gamma)\right\}
  +o(n^{-1}),
\end{equation}
again up to a factor of two depending on whether $f_M$ is defined as $F_M$ or
$F_M/2$.

The notation $Q_M$ is important. In the exact-fit, fixed-weight quadratic case,
$Q_M$ reduces to the familiar residual projector $U_M$ (up to the same factor
convention). Under fixed misspecification, estimated weights, or curved models,
$Q_M$ is the correct profile Hessian and $U_M$ is only a shortcut.

\section{Contrasts: Nested Tests, Vuong, and Absolute Fit}

Most SEM tests compare profiled values. Let
\[
  c(x)=\sum_{j=1}^J a_j v_{M_j}(x),
  \qquad
  s_c=\sum_j a_j s_{M_j},
  \qquad
  Q_c=\sum_j a_j Q_{M_j}.
\]
Then
\begin{equation}\label{eq:contrast-expansion}
  c(\hat x)-c(x_0)
  =
  n^{-1/2}s_c'Z
  +\frac{1}{2n}Z'Q_cZ
  +n^{-1}s_c'b_n
  +o_p(n^{-1}).
\end{equation}

\paragraph{Distinguishable regime.}
If $s_c\ne0$, the contrast is root-$n$ normal:
\[
  \sqrt n\{c(\hat x)-c(x_0)\}
  \dto
  N(0,s_c'\Gamma s_c).
\]
This is the regular Vuong-style regime and also the fixed-misspecification
regime for many fit-index differences.

\paragraph{Degenerate regime.}
If the null or comparison makes the values and first derivatives cancel, then
the quadratic term is leading:
\[
  n\,c(\hat x)
  \dto
  \frac12 Z'Q_cZ,
\]
with a weighted chi-square or indefinite quadratic-form reference distribution.
Nested difference tests are the main SEM example:
\[
  c(x)=v_B(x)-v_A(x),
  \qquad
  Q_c=Q_B-Q_A.
\]
Absolute goodness of fit is the special comparison with the saturated model:
\[
  v_{\mathrm{sat}}(x)\equiv0,
  \qquad
  Q_{\mathrm{sat}}=0,
  \qquad
  Q_c=Q_M.
\]
This explains why a single $U_M$ or $Q_M$ appears in RMSEA and exact-fit tests,
while nested tests naturally involve a difference $Q_B-Q_A$.

\paragraph{Local alternatives.}
Local alternatives belong here too. They perturb $c(x_0)$ or $s_c$ at the
$n^{-1}$ or $n^{-1/2}$ scale, turning central quadratic forms into noncentral
ones. This should be treated as a regular local regime, distinct from
nonregular boundary asymptotics.

\section{Fit Indices as Functionals}

Fit indices should be treated as smooth or nonsmooth functionals of the master
maps. For the first paper, the clean theorem is the smooth one. Let
\[
  y(x)=\{\theta_{M_1}(x),\ldots,\theta_{M_J}(x),
        r_{M_1}(x),\ldots,r_{M_J}(x),
        v_{M_1}(x),\ldots,v_{M_J}(x)\}
\]
and let $h$ be differentiable at $y(x_0)$. Then
\[
  \sqrt n\{h(\hat y)-h(y_0)\}
  \dto
  N(0,\nabla h'\,\Sigma_y\,\nabla h),
\]
where $\Sigma_y$ is assembled from the joint derivatives of the component maps
against the common $Z$.

This covers fixed-positive-misfit versions of:
\[
  \text{RMSEA, CFI, TLI, SRMR, information criteria, residual summaries,}
\]
provided denominators are nonzero and the functional is evaluated away from
kinks. Bias terms can be added by second-order delta method when useful.

The nonsmooth cases should be acknowledged rather than solved in this paper:
\[
  \max(\cdot,0),\quad \sqrt{\cdot}\text{ at zero},\quad clipped CFI/TLI,\quad
  near-zero baseline discrepancy,\quad boundary parameters.
\]
These are important, but they are not needed for a strong regular-asymptotics
paper.

\section{Implementation Roadmap in \texttt{magmaan}}

The implementation should mirror the theorem structure.

\begin{enumerate}
\item Define a first-stage object: $\hat x$, $x_0$ or fitted plug-in,
  $\Gamma$, optional $b$, and block metadata for moments, weights, groups, and
  ordinal pieces.
\item For each estimator/model, expose a profile object that can compute
  $f_M$, $\hat\theta_M$, $H_M$, $G_M$, residual derivatives, $s_M$, and $Q_M$.
\item Route standard errors through the coefficient theorem
  \eqref{eq:theta-vcov}.
\item Route residual diagnostics, score tests, modification indices, and SRMR
  ingredients through the residual-map theorem \eqref{eq:resid-derivative}.
\item Route exact-fit tests, nested tests, Vuong tests, RMSEA corrections, and
  future CFI/TLI corrections through the value-map theorem
  \eqref{eq:value-expansion}.
\item Keep $U_M$ as an optimized special case for exact-fit/fixed-weight
  quadratic models, but store and test the general object as $Q_M$.
\end{enumerate}

Validation should proceed from the simplest profile maps to the most coupled:
\[
\begin{array}{ll}
1. & \text{fixed-weight ULS/GLS, complete data;}\\
2. & \text{estimated-weight GLS/DWLS and ordinal limited information;}\\
3. & \text{ML/FIML as criterion-specific profile maps;}\\
4. & \text{multigroup and equality-constrained comparisons;}\\
5. & \text{fit-index functionals such as RMSEA, CFI, TLI, and SRMR.}
\end{array}
\]

\section{Proposed Paper Structure}

\begin{enumerate}
\item Introduction: widely used SEM outputs rely on exact-fit shortcuts; under
  misspecification, one needs profile derivatives at pseudo-true points.
\item Setup: first-stage CLT, smooth criteria, regular pseudo-true points.
\item Main theory: coefficient maps, residual maps, value maps.
\item Model comparisons: absolute fit, nested tests, Vuong and overlapping
  comparisons, local alternatives.
\item Fit indices: RMSEA now; CFI/TLI/SRMR as smooth functionals, with
  nonsmooth cases flagged.
\item SEM specializations: ML/FIML, GLS/ULS/WLS/DWLS, ordinal/multigroup.
\item Implementation: \texttt{magmaan} architecture, analytic derivatives,
  finite-difference validation, lavaan parity where relevant.
\item Simulations and examples: cases where exact-fit shortcuts agree, cases
  where $Q_M$ and observed-profile Hessians matter, and empirical size/coverage
  for robust SEs and tests.
\item Discussion: regular theory first; boundaries, rank loss, and clipped fit
  indices are future work.
\end{enumerate}

\section{Immediate Next Steps}

The next concrete work should be:
\begin{enumerate}
\item Write a short theorem note with consistent notation
  $(x,\Gamma,b,\theta_M,v_M,s_M,Q_M,H_M,G_M)$.
\item Refactor existing RMSEA and nested-test notes to cite that notation
  rather than each defining a parallel vocabulary.
\item Add a small fixed-weight validation script showing
  $n\{\E v_M(\hat x)-v_M(x_0)\}$ matches $\tfrac12\tr(Q_M\Gamma)$ under fixed
  misspecification.
\item Add an estimated-weight validation that isolates the weight channel.
\item Decide the first implemented public surface in \texttt{magmaan}: likely
  profile Hessians $Q_M$, robust nested tests, and fixed-misspecification RMSEA
  diagnostics before moving to CFI/TLI.
\end{enumerate}

\end{document}
